\section{Mechanized recovery and separating executions}
\label{sec:established}

\subsection{The multi-epoch symbolic theorem}

The checked symbolic model includes inherited roots, alternating proposers,
asymmetric confirmation, initial traffic, and unbounded epochs. Signing and
wrapping keys may be shared across their modeled uses. Its \textsf{Cut} event
exposes the selected wrapping keys; exposure of other epochs' entropy is
represented separately. We state the main theorem in terms of its recorded
events, preserving its original target and disclosure predicates.

Let $\tau$ be a trace of that model. Write
$\mathsf{SessionKeys}(s,a,b)$ for the actual identity keys associated with session
$s$, $\mathsf{KeyReserved}(s,e,P)$ and $\mathsf{EncapReserved}(s,e,Q)$ for the
target reservations, and $\mathsf{Message}(s,e,r,k,m)$ for an application-message
event. Define $\mathsf{AcquiredAfterSave}_c(s)$ to hold when there exist a role,
an epoch, a revision, and matching \textsf{Saved}/\textsf{Snapshot} events for
session $s$, with the \emph{save} event strictly after $c$. The matching record's
session projection is the relevant saved object.

\begin{theorem}[Checked multi-epoch symbolic secrecy]
\label{thm:symbolic}
Suppose $\tau$ contains the corresponding \textsf{SessionKeys}, \textsf{Cut}
at $c$, target reservation events at $g$ and $h$, and
$\mathsf{Message}(s,e,r,k,m)$. If
\[
 c<g,\qquad c<h,
\]
neither actual target identity key is disclosed anywhere in $\tau$, no reserved
entropy of target epoch $e$ is disclosed, and
$\mathsf{AcquiredAfterSave}_c(s)$ is false, then $\tau$ contains no adversarial
knowledge event $\mathsf K(m)$.
\end{theorem}

The theorem was checked with all 24 reusable dependencies in one Tamarin
process. An independent replay produced a byte-identical proof certificate.
Appendix~\ref{app:artifacts} identifies the exact query, source, toolchain,
replay receipt and complete auxiliary-query inventory.

The theorem uses the modeled cryptographic constructors and record interface.
Its key scope follows actual identity ownership: a shared signer exposed through
another session is still a disclosure of that target key. Likewise, a retained
target projection saved after Cut remains subject to the acquisition condition
when a sibling write carries it into a later aggregate image.

\subsection{Executability and the freshness boundary}

A finite two-round construction exercises the operational rules with an exposed
ancestor root and satisfies the complete original positive existential formula.
A separate source-derived evaluator checks its rule instances and global
quantifiers. This provides a concrete execution for the modeled recovery
premises.

The delayed-receipt construction instead contains a pre-Cut epoch-two
generation reservation and a post-Cut encapsulation. It derives the target
payload while satisfying the remaining main-theorem antecedents. Its point
audit identifies $c<g$ as the sole false antecedent. Seven countertraces are
rejected by the witness checker. Together, the positive and separating traces
connect the recovery condition to executable schedules.

\subsection{Native reservation disclosure}

The native experiment stops a child process after committing a reservation,
kills and reaps it, reopens the journal, and inspects the retained phase.
It exercises two distinct boundaries:
\begin{itemize}
\item a 277-byte key-generation token, before key computation or offer publication;
\item a 245-byte encapsulation token, after offer authentication and before
encapsulation computation or response publication.
\end{itemize}
The predictor obtains the previous root, journal identity, wrapping key and
captured token. Thereafter it uses the public policy configuration and later
authenticated control wires. Its prediction interface has no journal or
filesystem input, and the experiment withholds signer files and future private
owners. A wrong wrapping key is rejected.

For each cut, both honest peers complete the rekey and reopen their journals.
The driver then creates fresh 128-byte application payloads, sends three in each
direction, and checks honest delivery. Independently reconstructed traffic
owners authenticate and decrypt all six intercepted frames. The two cuts
therefore yield twelve actual post-restart plaintext recoveries. The
delayed-receipt extension repeats the six-message check at epoch two using one
responder snapshot.

\begin{proposition}[Execution-time and counter-based separation]
\label{prop:counter-separation}
In the native executions, computing a KEM after restart, completing
the target epoch, and subsequently reopening the journal do not make a
previously disclosed reservation an independent secret contribution.
Moreover, $\min(c_A,c_B)+1$ is not a universal upper bound on exposed epochs.
\end{proposition}
\begin{proof}
At both process-cut boundaries the captured reservation reconstructs the later
honest KEM result exactly. Equations~\eqref{eq:root}--\eqref{eq:traffic} then
reconstruct traffic keys from the captured predecessor root and public wires;
authenticated decryption verifies the six corresponding plaintexts after
completion and restart. In the delayed-receipt execution,
$(c_A,c_B)=(0,1)$ at capture and six epoch-two plaintexts are recovered.
Thus the proposed counter bound equals one while an exposed epoch equals two.
\end{proof}

\subsection{An aggregate write preserves the target origin}

A separate two-session native execution changes a sibling session and writes a
new aggregate image. It checks exact equality of the retained target record,
reopens and retries the original operation, and decrypts six subsequent target
messages in each build profile. The new global revision and outer encryption
nonce coexist with the old target reservation. This execution instantiates the
alias/provenance distinction in the record model.

The native executions isolate the reservation and acquisition conditions
used in Theorem~\ref{thm:symbolic}. Section~\ref{sec:computation} gives the
record simulations and common-experiment accounting for the corresponding
computational interfaces.

